
\section{Robustness and Extensions}
\label{sec:robustness}


Six groups of checks confirm the baseline result. (i)~FL-definition
variants: IQR multiplier, win-rate cutoff, temporal window.
(ii)~Sample and specification variants: unrestricted sample,
homogeneous subsamples, item$\times$year FE, matching.
(iii)~Alternative clustering. (iv)~Sensitivity to unobservables
(Cinelli--Hazlett). (v)~Staggered DiD.
Table~\ref{tab:robustness_summary} summarises all point estimates;
an additional column reports the identification assumption maintained
by each estimator.


\begin{table}[htbp]
\centering
\caption{Robustness summary: FL price coefficient across checks}
\label{tab:robustness_summary}
\small
\begin{tabular}{lcccll}
\hline\hline
Check & Coeff.\ & SE & $N$ & Assumption & Note \\
\hline
\multicolumn{6}{l}{\textit{Baseline and matching}} \\
OLS (item + year + PBU FE) & 0.064 & 0.020 & 1{,}654{,}401 & Sel.\ on obs.\ + FE & Baseline \\
CEM & 0.077 & 0.024 & 969{,}751 & Sel.\ on obs. & Sec.~\ref{sec:results_ols} \\
IPW & 0.055 & 0.021 & 830{,}194 & Sel.\ on obs. & Sec.~\ref{sec:results_ols} \\
Cross-fit average & 0.036 & --- & --- & Sel.\ on obs.\ + FE & No mechanical link \\
\hline
\multicolumn{6}{l}{\textit{FL definition variants}} \\
IQR 1.0$\times$ & 0.079 & 0.023 & 1{,}654{,}401 & Sel.\ on obs.\ + FE & 3{,}442 FL firms \\
IQR 1.5$\times$ (baseline) & 0.064 & 0.022 & 1{,}654{,}401 & Sel.\ on obs.\ + FE & 2{,}735 FL firms \\
IQR 2.0$\times$ & 0.060 & 0.022 & 1{,}654{,}401 & Sel.\ on obs.\ + FE & 2{,}093 FL firms \\
IQR 3.0$\times$ & 0.050 & 0.024 & 1{,}654{,}401 & Sel.\ on obs.\ + FE & 1{,}456 FL firms \\
\hline
\multicolumn{6}{l}{\textit{Specification variants}} \\
Item$\times$year FE & 0.074 & 0.022 & 1{,}654{,}401 & Sel.\ on obs.\ + FE & Tighter controls \\
Two-way clustering & 0.064 & 0.024 & 1{,}654{,}401 & Sel.\ on obs.\ + FE & Item + PBU \\
\hline
\multicolumn{6}{l}{\textit{Diagnostics (not primary estimates)}} \\
IV (leave-one-out) & 0.194 & 0.077 & 1{,}654{,}401 & Excl.\ restriction & Measurement error \\
\hline\hline
\end{tabular}
\begin{tablenotes}
\small
\item \textit{Notes:} ``Sel.\ on obs.\ + FE'' indicates that identification relies on selection on observables absorbed by item, year, and PBU fixed effects. Cross-fit further removes any mechanical link between the FL classification rule and the price outcome. The IV row is a measurement-error diagnostic; exclusion-restriction concerns preclude treating it as a primary estimate (Section~\ref{sec:results_ols}).
\end{tablenotes}
\end{table}


\paragraph{Threshold sensitivity.}
The FL coefficient decreases monotonically with stricter thresholds:
0.079 (IQR 1.0$\times$), 0.064 (baseline), 0.060 (2.0$\times$),
0.050 (3.0$\times$). This pattern likely reflects attenuation by
classification noise: stricter thresholds produce fewer FL firms and
fewer treated tenders, reducing statistical precision and increasing
the share of tenders misclassified as untreated. An alternative
interpretation---that marginal FL firms (near the threshold) operate
in markets with larger price effects while extreme-tail firms face
detection pressure that constrains cartel markups---cannot be excluded
but is less parsimonious.

\paragraph{Sensitivity to unobservables.}
\citet{cinelli2020making} sensitivity analysis yields $RV_{q=1} =
17.5\%$: an unobserved confounder would need to explain 17.5\% of
residual variation in \emph{both} FL presence and log prices to nullify
the coefficient---stringent given the item, year, and PBU fixed-effects
structure. Full sensitivity plots are in
Appendix~\ref{app:robustness}.


\paragraph{Staggered difference-in-differences.}
A \citet{callaway2021difference} staggered DiD yields a positive but
insignificant ATT of 0.014 (SE~$= 0.039$, MDE~$= 0.076$); the design
is underpowered for effects of the estimated magnitude. Full results,
including \citet{rambachan2023more} sensitivity analysis and
event-study plots, are in Appendix~\ref{app:did}.


\paragraph{Illustrative welfare estimates.}
Applying the price coefficients to FL-present procurement spending
as an order-of-magnitude exercise yields R\$734M (OLS, 2.4\% of BEC
spending) and R\$400M (cross-fit, 1.3\%) over 2009--2019. These
figures are upper bounds conditional on the causal interpretation of
the FL--price association; because our estimates capture a conditional
association, the true welfare impact may be smaller. Full derivation
and policy counterfactuals in Appendix~\ref{app:extensions}.


Oversight heterogeneity (Table~\ref{tab:regime_oversight}, Appendix~\ref{app:robustness}) shows FL price effects decreasing monotonically from small to large PBUs (Q1: 0.214 to Q4: 0.017), suggestive but underpowered. The threshold stability figure (Figure~\ref{fig:threshold_stability}) and IQR-vs-model-based classifier comparison are in Appendix~\ref{app:robustness}.
